
This thesis investigated how AI-generated structured DeFi intents
can be verified against authoritative execution relationships
before blockchain submission. The central problem addressed was
that individually valid protocols, contracts, functions, and assets
may still form an inconsistent execution structure. Neither a
well-formed language-model response nor wallet authorization alone
establishes that these entities belong together for a particular
operation. The objective was therefore to investigate a maintained,
domain-specific property graph as a pre-execution verification
mechanism and to examine how its decisions can be integrated
with controlled ERC-4337 wallet execution.

The proposed concept separates natural-language intent generation,
deterministic validation, authoritative knowledge acquisition,
action-specific relationship verification, and controlled execution.
The language model generates a candidate structured intent, while
execution-specific contracts, functions, and protocol resources
are resolved using admitted graph knowledge. This separation
prevents the language model from directly determining the
execution structure. The concept was implemented in a prototype
using Neo4j, covering selected Uniswap V3 single-hop swaps and
Aave V3 supply, withdrawal, borrowing, and repayment operations
on Ethereum Sepolia. Approved operations are retained server-side
and reverified before transaction preparation.

\medskip
\noindent\textbf{Contributions and Main Findings}

The principal contribution of this thesis is an action-specific
verification architecture that examines the relationships among
execution entities within the applicable blockchain network and
protocol deployment. The property graph represents maintained
verification knowledge against which structured operations are
checked. By explicitly evaluating the relationships required for
approval, the approach separates deterministic execution
verification from probabilistic language interpretation.
Furthermore, the verification mechanism distinguishes an
established incompatibility from insufficient evidence, ensuring
that neither condition can authorize execution.

The controlled relationship-mutation experiment provides the
strongest evidence for this contribution. Both the full graph
verifier and an entity-existence baseline approved all 12 valid
controls. After individual required relationships were removed,
the full verifier rejected all 60 altered execution paths,
whereas the baseline approved all 60. Because the individual
entities remained present, these results demonstrate the value
of verifying their required connections rather than checking
entity existence alone. The findings support the effectiveness
of the tested relationship patterns without establishing that
a graph database is the only suitable means of implementing
such verification.

Independent graph-verification experiments further supported
the proposed mechanism. All 340 reference-intent verification
calls, covering 112 distinct graph inputs, agreed with the
decisions and selected execution evidence derived from the
admitted graph snapshot. Of these calls, 338 (99.41\%) completed
within two seconds. Among the first occurrences of distinct
inputs, 110 of 112 calls (98.21\%) met the same threshold.
Both results satisfied the prototype-specific performance
requirement R6 within the recorded workload. In addition,
all 36 observations from the controlled knowledge-maintenance
experiment matched their predefined expectations, including
the prevention of approval when verification knowledge was
stale or incomplete.

The balanced natural-language evaluation produced 534 exact
final-status matches across 600 prompts, corresponding to
89.00\% agreement with assistant-drafted references. The more
restrictive joint-reference criterion was satisfied in 529
cases (88.17\%), while four attempts ended in provider quota
errors. These results indicate that the complete processing
pipeline remains less reliable than graph verification when
evaluated independently. They also demonstrate the importance
of separating intent-interpretation correctness from the
correctness of the resulting execution relationships.

The execution integration was initially evaluated through
application-level approval and handover controls. A subsequent
implementation extension introduced restricted ERC-4337
wallet authorization using a custom smart-account validator.
Local contract-level tests with the actual Kernel and EntryPoint
implementations demonstrated that execution authorization
could be bound to the complete approved transaction batch.
Unauthorized modifications to contracts, functions, parameters,
and authorization conditions were rejected. This extension
supported all four R4 assessment criteria within the locally
evaluated scope, while public Sepolia deployment and
bundler-level acceptance remained unverified.

Taken together, these findings demonstrate that the main
research objective was achieved within the admitted and
evaluated verification scope. The prototype establishes
action-specific execution relationships using maintained
graph knowledge and connects the resulting decisions to
controlled execution mechanisms.

\medskip
\noindent\textbf{Limitations}

The most significant limitation concerns the preservation
of the user's original request. In one evaluated case,
the language model omitted an explicitly requested
intermediate-token condition and generated a direct-swap
intent. The graph verifier approved this reduced operation
because its execution relationships were supported by the
admitted knowledge. No transaction was submitted. This
case demonstrates that correct relationship verification
does not necessarily establish correct interpretation of
the original request. A verifier may correctly assess the
structured operation it receives even when a user-specified
constraint has already been lost.

The evaluation also has methodological limitations.
The natural-language references were assistant-drafted
without documented independent item-level human review,
and the joint-reference measure examined only selected
aspects of the interpreted intents. Graph-verification
references were derived from the admitted snapshot rather
than independently confirming every underlying protocol
fact. In addition, the natural-language scenarios,
graph inputs, and mutation fixtures included dependent
observations. The earlier exploratory study and the
balanced follow-up used different datasets and procedures;
their results therefore cannot establish a causal
improvement in performance.

The reliability of verification also depends on the
correctness, freshness, and completeness of the underlying
knowledge. Although the synchronization procedures
automate knowledge acquisition when invoked, continuous
monitoring of authoritative sources was not implemented.
Manually maintained function mappings, non-atomic updates,
incomplete removal of obsolete relationships, and limited
representation of live reserve restrictions further
constrain the verification scope. Consequently, approval
establishes consistency with admitted graph knowledge,
not complete transaction feasibility under current
blockchain conditions.

Execution authorization and traceability have additional
limitations. The original wallet-handover evaluation
included 28 supplementary tests, several of which relied
on simulated outcomes. The subsequent restricted
authorization extension demonstrated wallet-level
enforcement using actual smart-account contracts in a
local EVM environment. However, the revised validator
was not installed on public Ethereum Sepolia, and its
compatibility with public bundlers and production
conditions was not established. The mechanism also
depends on a trusted verification authority, while
authenticated binding to the connected browser-wallet
owner and cumulative spending restrictions across
separate approvals remain incomplete. The custom
validator has not undergone an independent security
audit.

Traceability improved through an additional diagnostic
extension, which correctly localized all 60 controlled
relationship deletions and reported the corresponding
expected and observed conditions. Nevertheless, the
approved controls lacked an explicit reason or stable
verification-rule identifier, and some broader graph
results retained generic diagnostics. These findings
support improved traceability for the tested mutation
cases but do not establish complete R5 coverage across
all verification outcomes.

\medskip
\noindent\textbf{Future Work}

An important direction for future research is stronger
preservation of explicit user constraints throughout
intent generation, validation, verification, and
transaction preparation. This could include additional
mechanisms for checking consistency between the original
natural-language request and the generated structured
intent. Such improvements should be evaluated using
independently reviewed references and test requests
that are kept separate from development. Future
evaluations should also distinguish incorrect constraint
extraction from incorrect enforcement of constraints
that were extracted successfully.

Further work should improve the maintenance and
reliability of authoritative verification knowledge.
Potential extensions include continuous synchronization,
consistent and atomic graph updates, dependency-aware
freshness checks, and more comprehensive validation
of protocol-specific state. Expanding the evaluation
to additional protocols, blockchain networks, and
DeFi operations would help determine which relationship
patterns can be reused and which require specialized
verification rules. Beyond individual operations,
multi-step DeFi workflows present a further research
question: how can a sequence of individually valid
actions be verified consistently when the output
or state changes of one action influence the
requirements of subsequent actions?

The restricted wallet-authorization extension provides
a basis for further integration between graph verification
and blockchain execution. Future work should investigate
installation and operation on public test networks,
compatibility with ERC-4337 bundlers, independent security
auditing, authenticated wallet-owner authorization, and
cumulative spending policies. Additional evaluation is
also needed to examine how verification and execution
remain consistent when graph knowledge or blockchain
state changes between authorization and submission.

Finally, a systematic comparison with equally expressive
rule-based or relational verification implementations
could clarify the practical advantages and costs of
the property-graph representation. Such a comparison
would help distinguish the contribution of explicit
relationship verification from the particular database
technology used to implement it.

Overall, this thesis demonstrates that explicit,
action-specific relationship verification can provide
an additional safeguard between AI-generated DeFi
intents and blockchain execution. Its central contribution
is the establishment of a verifiable connection between
a structured operation and authoritative execution
knowledge. The effectiveness of this safeguard ultimately
depends on faithful interpretation of the user's
requirements, reliable maintenance of verification
knowledge, and enforcement of the approved operation
at the execution boundary.
